% chapter: wiki_sparam
\chapter{S-Parameters}\label{ch:wiki_sparam}

S-parameters (scattering parameters) describe how RF signals are scattered by a network. Unlike Z or Y parameters, S-parameters are measured with all ports terminated in a reference impedance Z₀ (typically 50 Ω), making them practical to measure with a vector network analyser (VNA) at microwave frequencies.

				

\section{Definition}

				

\rffig{wiki_sparam_1.pdf}{2-port S-parameter definitions. All ports terminated in Z₀ during measurement.}

				

For a 2-port network, S-parameters relate incident (a) and reflected/transmitted (b) wave amplitudes:

				
\[\begin{pmatrix}b_1\\b_2\end{pmatrix} = \begin{pmatrix}S_{11}&S_{12}\\S_{21}&S_{22}\end{pmatrix}\begin{pmatrix}a_1\\a_2\end{pmatrix}\]

				
\begin{itemize}

					\item S₁₁ — input reflection coefficient (return loss at port 1)

					\item S₂₁ — forward transmission (insertion loss or gain)

					\item S₁₂ — reverse isolation

					\item S₂₂ — output reflection coefficient

				
\end{itemize}

				

\section{Smith Chart}

				

\rffig{wiki_sparam_2.pdf}{Smith Chart: constant-resistance circles (purple), constant-reactance arcs (blue). Example load shown as orange dot.}

				

The Smith Chart maps the complex reflection coefficient Γ onto a plot of normalised impedance. Constant-resistance circles are in the horizontal centre; constant-reactance arcs run vertically. Moving clockwise along a constant-|Γ| circle corresponds to adding transmission line length toward the generator.

				

\section{Touchstone Files}

				

Measured S-parameters are stored in Touchstone format (.s1p, .s2p, etc.). Each line contains a frequency point followed by pairs of values encoding each Sij in MA (magnitude-angle), RI (real-imaginary), or DB (dB-angle) format.
